
\section{同伦幺元的代数框架}
\label{sec:3-3}

第 3.5 节将构造与拉格朗日子流形相伴的滤过 $A_\infty$ 代数。一般而言，
我们不知道这个滤过 $A_\infty$ 代数是否具有定义 3.2.20 (3) 意义下的幺元；
不过，第 7.3 节将说明它具有同伦幺元\liftIndexRangeStart{I099}\liftIndex{I102}。
本节给出同伦幺元的代数定义；与拉格朗日子流形相伴的 $A_\infty$ 代数之同伦幺元，
将在第 7.3 节中给出几何实现。第 3.3.2 小节还将定义保持幺元（相应地，保持同伦幺元）
的含幺（相应地，同伦含幺）$A_\infty$ 同态。

\subsection{同伦幺元的定义}
\label{subsec:3-3-1}

令 $(C,\mathfrak m)$ 为滤过 $A_\infty$ 代数。（我们主要取
$\Lambda_{0,\mathrm{nov}}$ 为系数环，但下面的框架显然也适用于 $R$ 上无滤过的
$A_\infty$ 代数。）这里不假设它含幺。

先准备一些记号。以 $B_k(B(C[1]))$ 表示 $k$ 个 $B(C[1])$ 的张量积，并置
\[
  B(B(C[1]))
  =\bigoplus_{k\geq1}B_k(B(C[1]))
  =B(C[1])\oplus
  \bigl(B(C[1])\otimes B(C[1])\bigr)\oplus\cdots .
\]
令 $\widehat B(\widehat B(C[1]))$ 为 $B(B(C[1]))$ 关于能量滤过的完备化。
为了区分 $B(C[1])$ 内部的张量积与 $B(B(C[1]))$ 中的张量积，我们把后者记作
$\mathbin{\overline\otimes}$，并用粗体字母 $\mathbf x$ 等表示 $B(C[1])$ 的元素。
例如，$x_1\otimes x_2$ 属于 $B_2(C[1])$，而
$\mathbf x_1\mathbin{\overline\otimes}\mathbf x_2$ 属于
$B_2(B(C[1]))$。若
\[
  \mathbf x_i=x_{i,1}\otimes\cdots\otimes x_{i,\ell_i}
  \in B_{\ell_i}(C[1]),
\]
则
\[
\begin{gathered}
  \mathbf x_1\mathbin{\overline\otimes}\cdots
  \mathbin{\overline\otimes}\mathbf x_k\in B_k(B(C[1])),\\
  \mathbf x_1\otimes\cdots\otimes\mathbf x_k
  \in B_{\ell_1+\cdots+\ell_k}(C[1]),
\end{gathered}
\]
而且
\[
  \mathbf x_1\otimes\cdots\otimes\mathbf x_k
  =x_{1,1}\otimes x_{1,2}\otimes\cdots\otimes x_{1,\ell_1}
  \otimes x_{2,1}\otimes\cdots\otimes x_{k,\ell_k}.
\]

现在置
\begin{equation}
  C^+=C\oplus\Lambda_{0,\mathrm{nov}}\mathbf e^+
  \oplus\Lambda_{0,\mathrm{nov}}\mathbf f,
  \tag{3.3.1}\label{eq:3-3-1}
\end{equation}
其中 $\deg\mathbf e^+=0$、$\deg\mathbf f=-1$；这里的次数是移位前次数。

\begin{definition}{3.3.2}\label{def:3-3-2}
令 $(C,\mathfrak m)$ 为 $A_\infty$ 代数。若元素
$\mathbf e\in C^0=C[1]^{-1}$ 具有如下性质，就称它是
$(C,\mathfrak m)$ 的\emph{同伦幺元}：存在保持能量滤过、次数为 $3-2k$ 的映射
\[
  \mathfrak h_k\colon B_k(B(C[1]))\longrightarrow C[1],
  \qquad k=1,2,\ldots,
\]
使得
\begin{equation}
  \mathfrak h_1=\mathfrak m,
  \tag{3.3.3}\label{eq:3-3-3}
\end{equation}
并使下面的 $(C^+,\mathfrak m^+)$ 成为定义 3.2.20 (3) 意义下以
$\mathbf e^+\in C^+[1]^{-1}$ 为幺元的 $A_\infty$ 代数。这里的 $C^+$ 如
\eqref{eq:3-3-1}，而 $\mathfrak m^+$ 定义如下：对 $x\in C^+$，
\begin{equation}
\begin{gathered}
  \mathfrak m_2^+(\mathbf e^+,x)
  =(-1)^{\deg x}\mathfrak m_2^+(x,\mathbf e^+)=x,\\
  \mathfrak m_k^+(\ldots,\mathbf e^+,\ldots)=0,
  \qquad k\ne2,
\end{gathered}
\tag{3.3.4}\label{eq:3-3-4}
\end{equation}
并且
\begin{equation}
\begin{aligned}
  &\mathfrak m^+
  (\mathbf x_1\otimes\mathbf f\otimes\mathbf x_2
  \otimes\mathbf f\otimes\cdots\otimes\mathbf f\otimes\mathbf x_k)\\
  &\qquad={}
  \mathfrak h_k(\mathbf x_1\mathbin{\overline\otimes}\cdots
  \mathbin{\overline\otimes}\mathbf x_k),
\end{aligned}
\tag{3.3.5.1}\label{eq:3-3-5-1}
\end{equation}
但 $k=2$ 且
$\mathbf x_1,\mathbf x_2\in\Lambda_{0,\mathrm{nov}}
(=\Lambda_{0,\mathrm{nov}}\cdot1)$ 时除外。在后一情形，置
\begin{equation}
  \mathfrak m_1^+(\mathbf f)
  =\mathbf e^+-\mathbf e
  +\mathfrak h_2(1\mathbin{\overline\otimes}1),
  \tag{3.3.5.2}\label{eq:3-3-5-2}
\end{equation}
并要求
\begin{equation}
  \mathfrak h_2(1\mathbin{\overline\otimes}1)
  \equiv0\pmod{\Lambda^+_{0,\mathrm{nov}}C^+}.
  \tag{3.3.5.3}\label{eq:3-3-5-3}
\end{equation}
还假定 $\mathfrak m_k^+$ 满足 (3.2.13)。换言之，假定
\[
  \mathfrak h_k\bmod\Lambda^+_{0,\mathrm{nov}}\colon
  B_k(B(\overline C[1]))\otimes_RR[e,e^{-1}]
  \longrightarrow\overline C[1]\otimes_RR[e,e^{-1}]
\]
由一个 $R$-模同态
\begin{equation}
  \overline{\mathfrak h}_k\colon
  B_k(B(\overline C[1]))\longrightarrow\overline C[1]
  \tag{3.3.6}\label{eq:3-3-6}
\end{equation}
诱导。
\end{definition}

称满足 (3.3.3)--(3.3.6) 的 $(C,\mathfrak m)$ 为\emph{带同伦幺元的滤过
$A_\infty$ 代数}，或简称\emph{同伦含幺滤过 $A_\infty$ 代数}
\liftIndexRangeEnd{I099}。

\begin{remark}{3.3.7}\label{rem:3-3-7}
\begin{enumerate}[label=\textup{(\arabic*)},leftmargin=2.8em]
  \item 若把 $\mathbin{\overline\otimes}$ 的次数计作
  $\deg'\mathbin{\overline\otimes}=-2$（即 $\mathbf f$ 的移位次数），
  则算子 $\mathfrak h_k$ 总是次数 $+1$ 的。

  \item 下面说明同伦幺元定义的几何来源。令 $(X,x_0)$ 为带基点空间，并假设存在
  乘法结构 $\mathord\cdot\colon X\times X\to X$，它在
  $\{x_0\}\times X$ 与 $X\times\{x_0\}$ 上的限制都同伦于恒等映射
  （换言之，$X$ 是 $H$-空间）。令
  \[
  \begin{aligned}
    H_1&\colon\{x_0\}\times X\times[0,1]\longrightarrow X,\\
    H_2&\colon X\times\{x_0\}\times[0,1]\longrightarrow X
  \end{aligned}
  \]
  分别为这样的同伦，即
  \[
  \begin{gathered}
    H_1(x_0,x,1)=H_2(x,x_0,1)=x,\\
    H_1(x_0,x,0)=x_0\mathbin{\cdot}x,\qquad
    H_2(x,x_0,0)=x\mathbin{\cdot}x_0.
  \end{gathered}
  \]
  那么 $x_0$ 是同伦幺元。取不交并 $X\sqcup[0,1]$，再把 $x_0$ 与 $0$
  识别，也就是置
  \[
    X_+=X\sqcup[0,1]/(x_0\sim0).
  \]
  把乘法延拓成 $\mathord\cdot\colon X_+\times X_+\to X_+$：若
  $s,s'\in[0,1]$ 且 $x,x'\in X$，则令
  \[
    s\mathbin{\cdot}x'=H_1(x_0,x',s),\qquad
    x\mathbin{\cdot}s'=H_2(x,x_0,s').
  \]
  因而 $1\cdot x=x\cdot1=x$，即 $1$ 是幺元。定义 3.3.2 中
  $(C^+,\mathfrak m^+)$ 的构造正是这个构造的代数版本。更准确地说，至少在经典层面，
  上面的观察给出了同伦幺元的几何含义。在这里，(3.3.5.1) 中含有
  $\mathfrak h_2(1\mathbin{\overline\otimes}1)$ 的公式，也就是
  (3.3.5.2)，并不出现。第 7.3 节将更详细地解释滤过 $A_\infty$ 代数中
  $\mathfrak h_2(1\mathbin{\overline\otimes}1)$ 一项的几何含义。

  \item 定义 3.3.2 中映射族 $\mathfrak h$ 的选择并不唯一，所以
  $(C^+,\mathfrak m^+)$ 不能由 $(C,\mathfrak m)$ 唯一确定。不过我们将证明，
  $(C^+,\mathfrak m^+)$ 作为滤过 $A_\infty$ 代数的同伦等价类
  （见第 4.2 节）由 $(C,\mathfrak m)$ 唯一确定。事实上，
  $(C^+,\mathfrak m^+)$ 与 $(C,\mathfrak m)$ 作为滤过 $A_\infty$ 代数同伦等价；
  见引理 4.2.55 和命题 4.2.52。在这个意义下，可以把 $C$ 看作类似于
  $C^+$ 的一个形变收缩核。
\end{enumerate}
\end{remark}

$(C^+,\mathfrak m^+)$ 的 $A_\infty$ 公式蕴含映射 $\mathfrak h_k$ 之间的若干关系。
现在把这些关系明确写出。定义
\[
\begin{aligned}
  \widehat{\mathfrak h}_k&\colon
  B(B(C[1]))\longrightarrow B(B(C[1])),\\
  \widehat{\mathfrak h}&\colon
  B(B(C[1]))\longrightarrow B(B(C[1]))
\end{aligned}
\]
如下。令
$\Delta\colon B(C[1])\to B(C[1])\otimes B(C[1])$ 为余积。对 $k\geq2$，置
\[
\begin{aligned}
 &\widehat{\mathfrak h}_k
 (\mathbf x_1\mathbin{\overline\otimes}\cdots
 \mathbin{\overline\otimes}\mathbf x_n)\\
 &={}\sum_{a,b,i}
 (-1)^{\deg'\mathbf x_1+\cdots+\deg'\mathbf x_{i-1}
 +\deg'\mathbf x'_{i,a}}\,
 \mathbf x_1\mathbin{\overline\otimes}\cdots
 \mathbin{\overline\otimes}\mathbf x_{i-1}\\
 &\quad{}\mathbin{\overline\otimes}
 \Bigl(
 \mathbf x'_{i,a}\otimes
 \mathfrak h_k(
 \mathbf x''_{i,a}\mathbin{\overline\otimes}\mathbf x_{i+1}
 \mathbin{\overline\otimes}\cdots
 \mathbin{\overline\otimes}\mathbf x_{i+k-2}
 \mathbin{\overline\otimes}\mathbf x'_{i+k-1,b})
 \otimes\mathbf x''_{i+k-1,b}
 \Bigr)\\
 &\quad{}\mathbin{\overline\otimes}\mathbf x_{i+k}
 \mathbin{\overline\otimes}\cdots
 \mathbin{\overline\otimes}\mathbf x_n,
\end{aligned}
\]
其中
\[
  \Delta\mathbf x_i
  =\sum_a\mathbf x'_{i,a}\otimes\mathbf x''_{i,a}.
\]
注意，$\widehat{\mathfrak h}_k$ 把 $B_n(B(C[1]))$ 映到
$B_{n-k+1}(B(C[1]))$；这是因为上面第三行括号内的
$\otimes$ 并不是 $\mathbin{\overline\otimes}$。这里 $\deg'$ 是移位后次数：
若 $\mathbf x=x_1\otimes\cdots\otimes x_m$，则
\[
  \deg'\mathbf x=\deg x_1+\cdots+\deg x_m-m;
\]
参见 (3.2.2)。对 $k=1$，置\translatornote{原书此式在求和号下写作
$a,i$，但外层被加数中没有自由指标 $a$；该指标只在紧接着展开
$\widehat d\mathbf x_i$ 时出现。译文删去外层游离的 $a$，保留对 $i$ 的求和。}
\[
\begin{aligned}
 &\widehat{\mathfrak h}_1
 (\mathbf x_1\mathbin{\overline\otimes}\cdots
 \mathbin{\overline\otimes}\mathbf x_n)\\
 &\quad={}
 \sum_i(-1)^{\deg'\mathbf x_1+\cdots+\deg'\mathbf x_{i-1}}\,
 \mathbf x_1\mathbin{\overline\otimes}\cdots
 \mathbin{\overline\otimes}\mathbf x_{i-1}
 \mathbin{\overline\otimes}\widehat d\mathbf x_i
 \mathbin{\overline\otimes}\mathbf x_{i+1}
 \mathbin{\overline\otimes}\cdots
 \mathbin{\overline\otimes}\mathbf x_n.
\end{aligned}
\]
这里回忆
\[
  \widehat d\mathbf x_i
  =\sum_a(-1)^{\deg'\mathbf x_{i,a}^{3;1}}\,
  \mathbf x_{i,a}^{3;1}\otimes
  \mathfrak m(\mathbf x_{i,a}^{3;2})\otimes
  \mathbf x_{i,a}^{3;3},
\]
其中使用了 (Conv.3) 的 (1.34)。

置 $\widehat{\mathfrak h}=\sum_k\widehat{\mathfrak h}_k$，并把它延拓到
$\widehat B(\widehat B(C[1]))$。对给定的
$\mathbf e\in C^0=C[1]^{-1}$，再引入映射族
\[
  \mathfrak b_{k,i}\colon B_k(B(C[1]))
  \longrightarrow B_{k-1}(B(C[1])),
  \qquad i=1,\ldots,k-1,
\]
定义为
\[
\begin{aligned}
 &\mathfrak b_{k,i}
 (\mathbf x_1\mathbin{\overline\otimes}\cdots
 \mathbin{\overline\otimes}\mathbf x_k)\\
 &\quad={}
 (-1)^{\deg'\mathbf x_1+\cdots+\deg'\mathbf x_i}\,
 \mathbf x_1\mathbin{\overline\otimes}\cdots
 \mathbin{\overline\otimes}
 (\mathbf x_i\otimes\mathbf e\otimes\mathbf x_{i+1})
 \mathbin{\overline\otimes}\cdots
 \mathbin{\overline\otimes}\mathbf x_k .
\end{aligned}
\]
若把 $\deg'\mathbin{\overline\otimes}$ 计作 $-2$，则
$\mathfrak b_{k,i}$ 的次数为 $+1$。

\begin{proposition}{3.3.8}\label{prop:3-3-8}
令 $(C,\mathfrak m)$ 为 $A_\infty$ 代数，则陈述 (3.3.9) 与 (3.3.10) 等价。
\end{proposition}

\par\smallskip\noindent\textup{(3.3.9)}\quad
定义 3.3.2 所定义的 $(C^+,\mathfrak m^+)$ 是以 $\mathbf e^+$ 为幺元的
$A_\infty$ 代数。

\par\smallskip\noindent\textup{(3.3.10)}\quad
映射 $\mathfrak h_k\colon B_k(B(C[1]))\to C[1]$ 满足
(3.3.5.2)、(3.3.5.3)、(3.3.6)，并具有如下性质：
\begin{equation}
  \mathfrak h_1=\mathfrak m.
  \tag{3.3.10.0}\label{eq:3-3-10-0}
\end{equation}
\begin{equation}
  \sum_{\ell=1}^{k}
  \mathfrak h_\ell\circ\widehat{\mathfrak h}_{k-\ell+1}
  -\sum_{i=1}^{k-1}\mathfrak h_{k-1}\circ\mathfrak b_{k,i}
  +\mathfrak P_k=0.
  \tag{3.3.10.1}\label{eq:3-3-10-1}
\end{equation}
这里当 $k\ne2$ 时 $\mathfrak P_k=0$；此外，在
\[
  \bigl(B_0(C[1])\mathbin{\overline\otimes}B_1(C[1])\bigr)
  \oplus
  \bigl(B_1(C[1])\mathbin{\overline\otimes}B_0(C[1])\bigr)
  \subset B_2(B(C[1]))
\]
的补集上，$\mathfrak P_2=0$；而在上述子空间上，
\begin{equation}
  -\mathfrak P_2(x\mathbin{\overline\otimes}1)
  =\mathfrak P_2(1\mathbin{\overline\otimes}x)=x.
  \tag{3.3.10.2}\label{eq:3-3-10-2}
\end{equation}

上面的 (3.3.10) 使我们无需引入 $(C^+,\mathfrak m^+)$，便可直接用原
$A_\infty$ 代数 $(C,\mathfrak m)$ 本身来表述其同伦幺元的定义。
命题 3.3.8 的证明只是直接计算，从略。

\subsection{含幺（相应地，同伦含幺）$A_\infty$ 同态}
\label{subsec:3-3-2}

本小节定义保持幺元或同伦幺元的 $A_\infty$ 同态。令
$(C_1,\mathfrak m^1)$、$(C_2,\mathfrak m^2)$ 为滤过 $A_\infty$ 代数，并令
$\mathfrak f\colon C_1\to C_2$ 为滤过 $A_\infty$ 同态。先假定 $C_i$ 以
$\mathbf e_i$ 为幺元（$i=1,2$）。

\begin{definition}{3.3.11}\label{def:3-3-11}
若
\begin{equation}
  \mathfrak f_k(x_1,\ldots,x_{i-1},\mathbf e_1,
  x_{i+1},\ldots,x_k)=0,\qquad k\geq2,
  \tag{3.3.12.1}\label{eq:3-3-12-1}
\end{equation}
并且
\begin{equation}
  \mathfrak f_1(\mathbf e_1)=\mathbf e_2,
  \tag{3.3.12.2}\label{eq:3-3-12-2}
\end{equation}
就称 $\mathfrak f$ \emph{保持幺元}，有时也称它为\emph{含幺
$A_\infty$ 同态}\liftIndex{I101}。
\end{definition}

无滤过版本以同样方式定义。

下面考虑具有同伦幺元的 $A_\infty$ 代数。假定滤过 $A_\infty$ 代数
$(C_1,\mathfrak m^1)$ 与 $(C_2,\mathfrak m^2)$ 分别具有同伦幺元
$\mathbf e_1$ 与 $\mathbf e_2$，并依 (3.3.1) 置
\translatornote{原书此式把最后一个生成元写成 $\mathbf f_i^+$，但它明言这里依
(3.3.1) 定义，而 (3.3.1) 的生成元是 $\mathbf f$；译文据此去掉不相容的上标。}
\[
  C_i^+=C_i\oplus\Lambda_{0,\mathrm{nov}}\mathbf e_i^+
  \oplus\Lambda_{0,\mathrm{nov}}\mathbf f_i.
\]

\begin{definition}{3.3.13}\label{def:3-3-13}
若滤过 $A_\infty$ 同态 $\mathfrak f\colon C_1\to C_2$ 可以延拓成含幺滤过
$A_\infty$ 同态 $\mathfrak f^+\colon C_1^+\to C_2^+$，就称
$\mathfrak f$ \emph{保持同伦幺元}，或称它为\emph{同伦含幺
$A_\infty$ 同态}。
\end{definition}

\section{杯积的 $A_\infty$ 形变}
\label{sec:3-4}

本节利用度量带状树空间与常值伪全纯圆盘的模空间，对每个定向但不必相对 spin 的
拉格朗日子流形，几何地构造与之相伴的经典 $A_\infty$ 代数
\liftIndex{I100}。（事实上，这一构造适用于任意定向流形 $L$。）
这里“经典”是指无滤过的 $A_\infty$ 代数；它定义在 $\mathbb Z$ 上。
应当指出，在经典 $A_\infty$ 代数的几何构造中，大对角线处的横截性问题起着关键作用。
相关证明将在第 7.2 节给出。本节只精确陈述所需的主要命题，并说明构造这一形变所需的
主要几何素材。

先回忆度量带状树的概念；更多细节见 \textbf{[FuOh97]}。在该文中，把度量带状树用于
莫尔斯理论和拓扑场论的想法首次得到严格实现。一个\emph{带状树}
\liftIndex{I103}是一个对子 $(T,i)$，其中 $T$ 是树，
$i\colon T\to D^2$ 是嵌入，并满足
\begin{align}
  &T\text{ 的任何顶点都不恰好连接两条边};
  \tag{3.4.1.1}\label{eq:3-4-1-1}\\
  &\text{若 }v\in T\text{ 是一价顶点，则 }i(v)\in\partial D^2;
  \tag{3.4.1.2}\label{eq:3-4-1-2}\\
  &i(T)\cap\partial D^2\text{ 由一价顶点组成}.
  \tag{3.4.1.3}\label{eq:3-4-1-3}
\end{align}
令 $G_k$ 为三元组 $(T,i,v_1)$ 的集合，其中 $(T,i)$ 是带状树的同痕类，
\[
  \#\bigl(i(T)\cap\partial D^2\bigr)=k,\qquad
  v_1\in i(T)\cap\partial D^2.
\]
位于 $\partial D^2$ 上的顶点称为\emph{外顶点}，其余称为\emph{内顶点}。
给外顶点排序，使 $v_1$ 为第一个，且该次序与 $\partial D^2$ 的逆时针定向相容。
含有外顶点的边称为\emph{外边}，其余称为\emph{内边}。

对每个 $\mathfrak t=(T,i,v_1)\in G_k$，分别以
$C_{\mathrm{ext}}^0(\mathfrak t)$、$C_{\mathrm{int}}^0(\mathfrak t)$ 表示
全体外顶点与全体内顶点之集；类似地，以
$C_{\mathrm{ext}}^1(\mathfrak t)$、$C_{\mathrm{int}}^1(\mathfrak t)$ 表示
全体外边与全体内边之集。

令 $\operatorname{Gr}(\mathfrak t)$ 为所有映射
\[
  \ell\colon C_{\mathrm{int}}^1(\mathfrak t)\longrightarrow\mathbb R^+
\]
组成的集合，并置
\[
  \operatorname{Gr}_k=\bigcup_{\mathfrak t\in G_k}
  \operatorname{Gr}(\mathfrak t).
\]
称其中的元素为\emph{度量带状树}：把 $\ell(e)$ 看作边 $e$ 的长度，并约定外边的
长度为无穷。$\operatorname{Gr}_k$ 上拓扑的描述见 \textbf{[FuOh97]}。
再以 $\overline{\operatorname{Gr}}(\mathfrak t)$ 表示所有映射
\[
  \ell\colon C_{\mathrm{int}}^1(\mathfrak t)
  \longrightarrow\mathbb R^+\cup\{\infty\}
\]
组成的集合，并把它们的并记作
$\overline{\operatorname{Gr}}_k$；这个空间是
$\operatorname{Gr}_k$ 的紧化。

下面依照 \textbf{[FuOh97]} 第 10 节定义映射
\[
  \Theta\colon\operatorname{Gr}_k
  \longrightarrow\mathcal M_k^{\mathrm{main,reg}}.
\]
取 $(\mathfrak t,\ell)\in\operatorname{Gr}_k$。对每条
$e\in C_{\mathrm{int}}^1(\mathfrak t)$，取欧氏矩形
$L_e=[0,\ell(e)]\times[0,1]$；对
$e\in C_{\mathrm{ext}}^1(\mathfrak t)$，取
$L_e=(-\infty,0]\times[0,1]$。分别从中删去
\[
  \{0,\ell(e)\}\times\{1/2\}
  \subset[0,\ell(e)]\times[0,1],
  \qquad
  \{0\}\times\{1/2\}
  \subset(-\infty,0]\times[0,1].
\]
给定 $v\in C_{\mathrm{int}}^0(\mathfrak t)$，考虑含有 $v$ 的两条相邻边
$e,e'$，其中 $e'$ 是依照含 $v$ 各边的循环次序紧接在 $e$ 后面的边。若
$e,e'$ 都是出边，就把
\[
  \{0\}\times(1/2,1]\subset\partial L_e
  \quad\text{与}\quad
  \{0\}\times[0,1/2)\subset\partial L_{e'}
\]
粘合；对 $e,e'$ 的其他定向情形也作类似粘合。见图 3.4.1。

\begin{figure}[htbp]
  \centering
  \setlength{\unitlength}{0.9mm}
  \begin{picture}(138,58)
    \put(27,55){\line(0,-1){9}}
    \put(37,55){\line(0,-1){9}}
    \put(27,46){\line(1,0){10}}
    \put(32,46){\circle{2}}
    \put(5,42){\line(1,0){20}}
    \put(25,42){\line(0,-1){10}}
    \put(5,32){\line(1,0){20}}
    \put(25,37){\circle{2}}
    \put(39,42){\line(1,0){20}}
    \put(39,42){\line(0,-1){10}}
    \put(39,32){\line(1,0){20}}
    \put(39,37){\circle{2}}
    \put(27,29){\line(1,0){10}}
    \put(27,29){\line(0,-1){13}}
    \put(37,29){\line(0,-1){13}}
    \put(27,16){\line(1,0){10}}
    \put(32,29){\circle{2}}
    \put(32,16){\circle{2}}
    \put(8,0){\line(4,3){18}}
    \put(15,0){\line(4,3){14}}
    \put(26,13.5){\circle{2}}
    \put(29,10.5){\circle{2}}
    \put(56,0){\line(-4,3){18}}
    \put(49,0){\line(-4,3){14}}
    \put(38,13.5){\circle{2}}
    \put(35,10.5){\circle{2}}
    \put(67,26){$\Longrightarrow$}
    \put(102,55){\line(0,-1){12}}
    \put(112,55){\line(0,-1){12}}
    \put(86,43){\line(1,0){16}}
    \put(112,43){\line(1,0){21}}
    \put(102,43){\line(1,-1){10}}
    \put(112,43){\line(-1,-1){10}}
    \put(83,33){\line(1,0){19}}
    \put(112,33){\line(1,0){24}}
    \put(102,33){\line(0,-1){14}}
    \put(112,33){\line(0,-1){14}}
    \put(102,19){\line(-3,-2){21}}
    \put(102,19){\line(1,-1){5}}
    \put(112,19){\line(3,-2){21}}
    \put(112,19){\line(-1,-1){5}}
    \put(107,14){\line(-3,-2){18}}
    \put(107,14){\line(3,-2){18}}
  \end{picture}
  \par\smallskip
  \renewcommand{\thefigure}{3.4.1}
  \refstepcounter{figure}
  \textbf{\figurename~\thefigure}\label{fig:3-4-1}
\end{figure}

这样便得到带有一个不完备平坦度量的空间 $X_0(\mathfrak t,\ell)$。
把 $X_0(\mathfrak t,\ell)$ 中的孔共形填上，得到带复结构和奇异黎曼度量的空间
$X(\mathfrak t,\ell)$。以
\[
  \Theta(\mathfrak t,\ell):=[X(\mathfrak t,\ell)]
\]
表示 $X(\mathfrak t,\ell)$ 的同构类。本书前两位署名作者在
\textbf{[FuOh97]} 的定理 10.4 中证明了
\[
  \Theta\colon\operatorname{Gr}_k
  \longrightarrow\mathcal M_k^{\mathrm{main,reg}}
\]
是同胚。（回忆 $\mathcal M_k^{\mathrm{main,reg}}$ 微分同胚于
$\mathbb R^{k-3}$；见第 2.1 节。）我们正是利用这个同胚在
$\operatorname{Gr}_k$ 上定义光滑结构。

映射 $\Theta$ 可以延拓到 $\operatorname{Gr}_k$ 的紧化，并给出同胚
\[
  \overline\Theta\colon\overline{\operatorname{Gr}}_k
  \longrightarrow\mathcal M_k^{\mathrm{main}}.
\]
由此，我们在 $\mathcal M_k^{\mathrm{main,reg}}$ 的每个元素上选定了一个奇异度量。
在奇点邻域内修改这族度量，就得到一族光滑的光滑黎曼度量。第 7.1 节将用这些度量
定义适当的索伯列夫范数。

现在为辛流形 $M$ 中的每个定向拉格朗日子流形 $L$ 构造相伴的经典
$A_\infty$ 代数。令
\[
  \Delta_{k+1}\subset L^{k+1}
  =\underbrace{L\times\cdots\times L}_{k+1\text{ 个}}
\]
为大对角线，故 $\Delta_{k+1}\cong L$；并令
$N_{\Delta_{k+1}}L^{k+1}$ 为 $\Delta_{k+1}$ 在 $L^{k+1}$ 中的法丛。
把这个法丛同 $\Delta_{k+1}$ 在 $L^{k+1}$ 中的一个管状邻域
$\operatorname{Tube}_{\Delta_{k+1}}L^{k+1}$ 识别。对
$\beta\in\Pi(M,L)$，以 $\mathcal M_{k+1}^{\mathrm{main}}(\beta)$ 表示如下
模空间的主分支：从 $(\Sigma,\partial\Sigma)$ 到 $(M,L)$、代表类 $\beta$、
边界上带 $k+1$ 个标记点的亏格零稳定伪全纯映射（见定义 2.1.27）。

在
\[
  \beta_0=0,\qquad k+1>2
\]
的情形，$\mathcal M_{k+1}^{\mathrm{main}}(\beta_0)$ 是常值映射的模空间。
这个模空间是横截的，即具有正确的维数 $k-2+n$。然而，评价映射
\[
  \operatorname{ev}\colon
  \mathcal M_{k+1}^{\mathrm{main}}(\beta_0)\longrightarrow L^{k+1}
\]
不是浸没。因此，在它上面取如下仓西结构（仓西结构的定义见第 A1 节）：
\[
  U_{k+1}
  =\operatorname{Tube}_{\Delta_{k+1}}L^{k+1}
  \times\overline{\operatorname{Gr}}_{k+1}.
\]
把法丛 $N_{\Delta_{k+1}}L^{k+1}$ 拉回到 $U_{k+1}$，以所得丛为障碍丛
$E_{k+1}$。利用识别
\[
  \operatorname{Tube}_{\Delta_{k+1}}L^{k+1}
  \cong N_{\Delta_{k+1}}L^{k+1},
\]
可以找到截面
$\mathfrak s_{k+1}\colon U_{k+1}\to E_{k+1}$，它恰好在
$\Delta_{k+1}\times\overline{\operatorname{Gr}}_{k+1}$ 上为零。
于是 $(U_{k+1},E_{k+1},\mathfrak s_{k+1})$ 定义一个仓西结构，其自同构群平凡；
而
\[
  \operatorname{ev}=(\operatorname{ev}_1,\ldots,\operatorname{ev}_k,
  \operatorname{ev}_0)\colon U_{k+1}\longrightarrow L^{k+1}
\]
是浸没（事实上，它是一个开嵌入与投影的复合）。这就是说，相对于这个仓西结构，
$\operatorname{ev}$ 是弱浸没的（定义 A1.13）。

为了构造滤过 $A_\infty$ 代数，还必须在
$\mathcal M_{k+1}^{\mathrm{main}}(\beta)$ 上选取适当的仓西结构；原因正是上面提到的
横截性问题。第 7 章的第 7.1、7.2 节将详细讨论这一点。在现阶段，本节暂不展开仓西
结构的讨论，只陈述下面的命题。

\begin{proposition}{3.4.2}\label{prop:3-4-2}
模空间 $\mathcal M_{k+1}^{\mathrm{main}}(\beta)$ 上存在一个定向仓西结构，
它与紧化相容；而在
$\mathcal M_{k+1}^{\mathrm{main}}(\beta_0)$ 上，该仓西结构与上面构造的仓西结构
一致。
\end{proposition}

命题 3.4.2 将在第 7.1 节证明，其中还会给出命题所说“相容”的精确定义
（命题 7.1.2）。粗略地说，含义如下：
$\mathcal M_{k+1}^{\mathrm{main}}(\beta)$ 的紧化中出现的边界分支，可以同若干
$\mathcal M_{k'+1}^{\mathrm{main}}(\beta')$ 的纤维积识别。带仓西结构空间的纤维积
具有诱导仓西结构（见第 A1.2 节）；相容性是说，由纤维积在边界上诱导的仓西结构，
与从 $\mathcal M_{k+1}^{\mathrm{main}}(\beta)$ 的仓西结构作为边界所诱导的仓西结构
一致。

下面陈述原仓西映射 $\mathfrak s$ 存在适当扰动这一主要结果。考虑 $L$ 上光滑奇异
单形的一个可数集 $\mathcal X_L$，要求只要 $P\in\mathcal X_L$，它的每个面
$\partial_iP$ 也属于 $\mathcal X_L$。之所以取可数集，是因为我们采用“光滑对应”
方法，并需要使用贝尔纲定理来作横截性论证；第 7 章将详细讨论这一点。把
$\mathcal X_L$ 中的元素写成 $P=(|P|,f)$，其中 $|P|$ 是单形，
$f\colon|P|\to L$ 是光滑映射。（有时滥用记号，以 $P$ 代替 $|P|$。）

考虑由 $\mathcal X_L$ 生成的自由 $\mathbb Z$-模，记作
$\mathbb Z\mathcal X_L$。我们将在 $\mathbb Z\mathcal X_L$ 上定义
$A_\infty$ 代数结构。令 $(P_i,f_i)\in\mathcal X_L$，并置
\begin{equation}
\begin{aligned}
 &\mathcal M_{k+1}^{\mathrm{main}}
   (\beta_0;P_1,\ldots,P_k)\\
 &\quad={}
 \mathcal M_{k+1}^{\mathrm{main}}(\beta_0)
 {}_{(\operatorname{ev}_1,\ldots,\operatorname{ev}_k)}
 \mathop{\times}_{f_1\times\cdots\times f_k}
 (P_1\times\cdots\times P_k).
\end{aligned}
\tag{3.4.3}\label{eq:3-4-3}
\end{equation}
式 \eqref{eq:3-4-3} 右端是仓西结构空间意义下的纤维积；一般定义见
第 A1.2 节的引理 A1.39。这里说明它在当前情形中的含义。我们已经在
$\mathcal M_{k+1}^{\mathrm{main}}(\beta_0)$ 上定义了仓西结构
$(U_{k+1},E_{k+1},\mathfrak s_{k+1})$。评价映射
\[
  \operatorname{ev}=(\operatorname{ev}_1,\ldots,\operatorname{ev}_k,
  \operatorname{ev}_0)\colon U_{k+1}\longrightarrow L^{k+1}
\]
是浸没，因而纤维积
\begin{equation}
\begin{aligned}
 U_{k+1}(P_1,\ldots,P_k)
 &={}
 U_{k+1}
 {}_{(\operatorname{ev}_1,\ldots,\operatorname{ev}_k)}
 \mathop{\times}_{f_1\times\cdots\times f_k}
 (P_1\times\cdots\times P_k)
\end{aligned}
\tag{3.4.4}\label{eq:3-4-4}
\end{equation}
是横截的，并且是一个带有光滑三角剖分的流形（可能带角）。
把丛 $E_{k+1}\to U_{k+1}$ 拉回，所得丛记作
$E_{k+1}(P_1,\ldots,P_k)$。$E_{k+1}$ 的截面 $\mathfrak s_{k+1}$
（仓西映射）在这个拉回丛上诱导一个截面，记作
$\mathfrak s_{P_1,\ldots,P_k}$。于是
\[
  \bigl(U_{k+1}(P_1,\ldots,P_k),
  E_{k+1}(P_1,\ldots,P_k),
  \mathfrak s_{P_1,\ldots,P_k}\bigr)
\]
在把 \eqref{eq:3-4-3} 右端纤维积视为拓扑空间时，为其定义了仓西结构。

从现在起，把 $\mathcal M_{k+1}^{\mathrm{main}}
(\beta_0;P_1,\ldots,P_k)$ 视为带仓西结构的空间；并且定义了强连续映射
\[
  \operatorname{ev}_0\colon
  \mathcal M_{k+1}^{\mathrm{main}}(\beta_0;P_1,\ldots,P_k)
  \longrightarrow L
\]
（定义见定义 A1.13）。

下面的命题将在第 7.2 节证明（命题 7.2.35）。置
$\vec P=(P_1,\ldots,P_k)$，其中 $P_i\in\mathcal X_L$。

\begin{proposition}{3.4.5}\label{prop:3-4-5}
可以选取 $\mathcal X_L$，并在
$\mathcal M_{k+1}^{\mathrm{main}}(\beta_0;\vec P)$ 上选取一套扰动
$\mathfrak s_{\vec P}$（即障碍丛的单值截面），使得下列性质成立：
\translatornote{原书 (3.4.6.4) 把包含关系与同构目标分别写成
$\mathbb Z\mathcal X_L\subset C(L;\mathbb Q)$ 和 $H(L;\mathbb Q)$。
但 $C(L;\mathbb Q)$ 到第 3.5 节才定义，而本节已将
$\mathbb Z\mathcal X_L$ 定义为 $S(L;\mathbb Z)$ 的子模；更直接地，
定理 3.4.8 的 (3.4.9.3) 明确给出 $\mathbb Z$ 系数结论并说它由
(3.4.6.4) 推出。译文据此统一为整数系数。}
\begin{align}
  &\mathfrak s_{\vec P}\text{ 横截于 }0;
  \tag{3.4.6.1}\label{eq:3-4-6-1}\\
  &\mathfrak s_{\vec P}\text{ 在紧化处相容};
  \tag{3.4.6.2}\label{eq:3-4-6-2}\\
  &\mathfrak s_{\vec P}^{-1}(0)\text{ 存在一个光滑三角剖分，使}\notag\\[-0.3em]
  &\quad\bigl(\mathfrak s_{\vec P}^{-1}(0),\operatorname{ev}_0\bigr)
  \text{ 可视为 }\mathcal X_L\text{ 中元素的}\notag\\[-0.3em]
  &\quad\mathbb Z\text{-线性组合};
  \tag{3.4.6.3}\label{eq:3-4-6-3}\\
  &\mathbb Z\mathcal X_L\subset S(L;\mathbb Z),\qquad
  \text{其包含映射诱导同构}\notag\\[-0.3em]
  &H(\mathbb Z\mathcal X_L,\delta)\cong H(L;\mathbb Z).
  \tag{3.4.6.4}\label{eq:3-4-6-4}
\end{align}
\end{proposition}

对 (3.4.6.3) 作些说明。由 (3.4.6.1)，
$\mathfrak s_{\vec P}^{-1}(0)$ 是光滑流形。给它取一个光滑三角剖分，并给各顶点
排序。利用顶点次序，可以把这个剖分的每个单形典范地同标准单形识别。把
\[
  \operatorname{ev}_0\colon\mathfrak s_{\vec P}^{-1}(0)\longrightarrow L
\]
限制到其中一个单形上，就得到 $L$ 的一个光滑奇异单形；(3.4.6.3) 断言它属于
$\mathcal X_L$。

令 $\Delta_a^d$（$a\in A$）为 $\mathfrak s_{\vec P}^{-1}(0)$ 中所有维数为
\[
  d=\dim\mathfrak s_{\vec P}^{-1}(0)
\]
的单形。以后记
\[
\begin{aligned}
 &\mathcal M_{k+1}^{\mathrm{main}}
   (\beta_0;P_1,\ldots,P_k)^{\mathfrak s}\\
 &\quad:={}
 \bigl(\mathfrak s_{\vec P}^{-1}(0),\operatorname{ev}_0\bigr)
 =\sum_{a\in A}\pm(\Delta_a^d,\operatorname{ev}_0),
\end{aligned}
\]
并把它看作 $L$ 的光滑奇异链。上式中正负号的定义如下：
模空间 $\mathcal M_{k+1}^{\mathrm{main}}(\beta_0;\vec P)$ 带有定向
（见第 8 章）；每个单形 $\Delta_a^d$ 都同带标准定向的
$\mathbb R^d$ 中的标准单形识别。当且仅当这两个定向一致时取正号
（见第 7.2 节注 7.2.15 (3)）。

现在定义
$\overline{\mathfrak m}^{\,L}
=\{\overline{\mathfrak m}_k^{\,L}\}_{k=0}^{\infty}$；它们将在
$\mathbb Z\mathcal X_L$ 上给出 $A_\infty$ 代数结构：
\begin{align}
  &\overline{\mathfrak m}_0^{\,L}=0,
  \tag{3.4.7.1}\label{eq:3-4-7-1}\\
  &\overline{\mathfrak m}_1^{\,L}(P)=(-1)^n\partial P,
  \tag{3.4.7.2}\label{eq:3-4-7-2}\\
  &\overline{\mathfrak m}_k^{\,L}(P_1,\ldots,P_k)
  =\mathcal M_{k+1}^{\mathrm{main}}
  (\beta_0;P_1,\ldots,P_k)^{\mathfrak s},
  \tag{3.4.7.3}\label{eq:3-4-7-3}
\end{align}
其中 $n=\dim L$。(3.4.7.2) 中的 $\partial$ 是通常的经典边界算子；
这里把 $P$ 看作一条链，另见注 3.5.8 (1)。

注意，命题 3.4.2 和 $\mathcal X_L$ 的选取保证了
(3.4.7.3) 右端的纤维积可以在仓西结构的意义下定义。
(3.4.7.3) 中模空间的定向在第 8 章定义 8.4.1 中给出。这里，经同构
\[
  \overline\Theta\colon
  \overline{\operatorname{Gr}}_{k+1}
  \longrightarrow\mathcal M_{k+1}^{\mathrm{main}}
  \cong\mathcal M_{k+1}^{\mathrm{main}}(\beta_0),
\]
把常值映射模空间 $\mathcal M_{k+1}^{\mathrm{main}}(\beta_0)$ 上的定向转移到
$\Delta_k\times\overline{\operatorname{Gr}}_{k+1}$。前一模空间的定向将在
第 8 章第 8.6 节给出。另一方面，第 2.1.2 小节已经在
$\mathcal M_{k+1}^{\mathrm{main,reg}}$ 上定义了另一个定向。分析
$\mathcal M_{k+1}^{\mathrm{main}}$ 上的定向并使用
\[
  \overline\Theta\colon\overline{\operatorname{Gr}}_{k+1}
  \longrightarrow\mathcal M_{k+1}^{\mathrm{main}},
\]
也可在 $\Delta_k\times\overline{\operatorname{Gr}}_{k+1}$ 上得到一个定向。
这个定向可能不同于上面把它视为零点集
$\mathfrak s_{k+1}^{-1}(0)$ 时所定义的定向。两者的差异可以精确算出，但因后文
不需要，这里从略。

由此得到下面的定理。

\begin{theorem}{3.4.8}\label{thm:3-4-8}
令 $L$ 为光滑定向流形，则 $L$ 上存在一个由光滑奇异单形组成的可数集
$\mathcal X_L$，满足下列性质：
\begin{align}
  &(\mathbb Z\mathcal X_L,\overline{\mathfrak m}^{\,L})
  \text{ 是 }\mathbb Z\text{ 上的 }A_\infty\text{ 代数};
  \tag{3.4.9.1}\label{eq:3-4-9-1}\\
  &\text{$L$ 存在一个光滑单纯剖分，使得可将 }
  PD[L]\text{ 看作 }\mathbb Z\mathcal X_L\text{ 的元素};\notag\\[-0.2em]
  &\text{此时 }PD[L]\text{ 是 }
  (\mathbb Z\mathcal X_L,\overline{\mathfrak m}^{\,L})\text{ 的同伦幺元};
  \tag{3.4.9.2}\label{eq:3-4-9-2}\\
  &H(\mathbb Z\mathcal X_L,\overline{\mathfrak m}_1^{\,L})
  \cong H(L;\mathbb Z).
  \tag{3.4.9.3}\label{eq:3-4-9-3}
\end{align}
这里 $PD$ 表示庞加莱对偶。
\end{theorem}

(3.4.9.1) 由命题 3.4.5，特别是 (3.4.6.2) 推出；详细证明见
第 7.2.4--7.2.8 小节。下面解释 (3.4.9.2)。给 $L$ 取一个光滑单纯剖分，并给
它的顶点排序。于是每个单形都可同标准单形识别，从而确定一个奇异单形。所有维数为
$\dim L$ 的这些奇异单形之和
\begin{equation}
  \sum_a\pm\Delta_a^{\dim L}
  \tag{3.4.10}\label{eq:3-4-10}
\end{equation}
代表基本同调类 $[L]$；这里的正负号由标准单形的定向是否同 $L$ 的定向一致而定
（见第 7.3.2 小节）。

本书把奇异链复形视作上链复形，因此 \eqref{eq:3-4-10} 被看作
$S^0(L;\mathbb Z)$ 的元素，并记作 $PD[L]$。(3.4.9.2) 将在第 7.3 节证明，
而 (3.4.9.3) 由 (3.4.6.4) 推出。称
$(\mathbb Z\mathcal X_L,\overline{\mathfrak m}^{\,L})$ 为\emph{与 $L$ 相伴的经典
$A_\infty$ 代数}。

文献 \textbf{[McC04, Wil05]} 中还有一些与定理 3.4.8 密切相关的工作。

\begin{remark}{3.4.11}\label{rem:3-4-11}
第 7.2 节并不直接使用扰动截面 $\mathfrak s_{\vec P}$ 构造
$A_\infty$ 结构。相反，对每个 $K$，我们先直接用
$\mathfrak s_{\vec P}$ 构造一个 $A_K$ 结构（见定义 4.4.6），然后得到一个与它
$A_K$-同伦等价的 $A_\infty$ 结构（定义 4.4.8）。不能一次直接构造
$A_\infty$ 结构的原因将在第 7.2.3 小节解释；这是一个相当技术性的要点。
本章其他构造同样有这个问题，但以后不再重复说明。实际使用时，一次构造
$A_\infty$ 结构，或先对每个 $K$ 构造 $A_K$ 结构、再以较间接的方式得到
$A_\infty$ 结构，并无实质区别。
\end{remark}

第 4 章将证明（作为该章定理 4.1.1 的一部分），$A_\infty$ 代数
$(\mathbb Z\mathcal X_L,\overline{\mathfrak m}^{\,L})$ 的弱同伦等价类由 $L$ 唯一确定。
也就是说，在弱同伦等价意义下，它与命题 3.4.5 所构造的扰动截面
$\mathfrak s^k$ 及可数集 $\mathcal X_L$ 的选择无关。
\translatornote{原书此处写“定理 3.4.8 所构造的扰动截面”，但扰动截面的存在与
性质陈述在命题 3.4.5；定理 3.4.8 陈述由其导出的 $A_\infty$ 代数。译文改正交叉引用。}
此外，第 7.5 节将证明：若置
\[
  \mathbb R\mathcal X_L=\mathbb Z\mathcal X_L\otimes\mathbb R
\]
并把 $\overline{\mathfrak m}^{\,L}$ 延拓到 $\mathbb R\mathcal X_L$，则
$(\mathbb R\mathcal X_L,\overline{\mathfrak m}^{\,L})$ 与 $L$ 在 $\mathbb R$ 上的德拉姆复形弱同伦等价
（定理 X）。
